% =====================================================================
%  Section 6: Conclusion
% =====================================================================
\section{Conclusion}
\label{sec:conclusion}

We presented \textsc{TADS} (Trajectory-Anchored Dynamic Selection),
a dynamic instruction-tuning data-selection framework that augments
an RL composite-reward base selector with a structural capability
anchor extracted from the model's own hidden-state trajectory. The
anchor is model-intrinsic---requiring no in-domain seed, labels,
validation set, or auxiliary scorer---and enters the selection score
as a single multiplicative term controlled by one hyperparameter,
with $\lambda{=}0$ recovering the base selector exactly. We proved a
stability guarantee (Theorem~\ref{thm:stability}): under a positive
eigenvalue gap, the trajectory anchor changes by at most a
learning-rate-scaled amount per epoch, and this bound transfers to
the selection criterion. To our knowledge, this is the first
stability guarantee for a dynamic IT data-selection method, and it
is base-agnostic rather than tied to a specific reward.

On the Qwen2.5-14B + Alpaca-GPT4 setting at $50\%$ selection,
\textsc{TADS} attains the best knowledge accuracy among the
evaluated baselines without per-task seed construction, while adding
under $0.5\%$ wall-clock overhead. A broader empirical study---across
model scales from 0.5B to 14B, multiple instruction-tuning datasets,
and an eight-task benchmark suite, together with diagnostic
experiments on baseline weaknesses---is in progress and described in
App.~\ref{app:inprogress}; it is designed to test whether the gains
reported here hold consistently across models, datasets, and
evaluation regimes under a single \textsc{TADS} configuration.